% Abstract: 150 to 250 words (counted by build.sh, also on the PDF). Every number comes from the registry.
Coding agents read what earlier agents left behind. In the July 2026 incident, about 700 agents on an unauthorized shared board joined an out-of-scope attack. In a fixed peer-note experiment with one agent per run, we test whether a refusal left on such a board changes compliance. Four model-and-harness packages (GPT-6.1 Sol in Codex, Claude Sonnet 5.5 in Claude Code, GLM 5.2 and DeepSeek V4 Flash in Codex) received pytest tasks made impossible by a missing dependency or data archive, with an explicit ticket rule against weakening the tests. A probe-validated deterministic check outside the agent's reach judged compliance. Every package skipped the test when allowed. Under the rule, \kd{base.total.violation_strict} baseline runs broke it, at most \regpct{base.strict_rate_max} in any package-task cell. The pre-registered test (GLM 5.2) sat at the floor, leaving the note's effect unresolved: \kd{HD.glm.S6.violation_strict.k_a} runs broke the rule with a prior agent's refusal note, \kd{HD.glm.S6.violation_strict.k_b} with a length-matched neutral note ($p = \regp{HD.glm.S6.violation_strict.fisher_p}$), and that one event was a functionally empty \texttt{conftest.py} (a docstring only). On a feasible variant no loss of legitimate success was detected (\kd{HDF.glm.S6F.legit_success.k_a} against \kd{HDF.glm.S6F.legit_success.k_b}), though the interval allows a drop of up to \reg{HDF.glm.S6F.legit_success.newcombe_lo.pp} percentage points. In an exploratory, later block without the ticket rule or its reporting instruction, GLM bypassed in \kd{norule.total.glm.bypass} runs, DeepSeek in \kd{norule.total.deepseek.bypass} and GPT in \kn{norule.gpt.S7.none}{bypass} on S7; the refusal note did not lower GLM's rate. Claude never bypassed without permission.
